\section{A uniform orbit criterion for quantitative width}\label{sec:uniform57}
The spherical estimate has an extension in which the dimension of the ambient feature space is not the relevant exponent. Two geometric quantities enter: a uniform small-ball exponent for every possible conditional-centroid direction, and the covering dimension of the particular target orbit. They need not agree. Keeping this distinction is essential for matrix-valued experiments.

Let $G$ be a compact connected Lie group with normalized Haar measure $m_G$ and a bi-invariant Riemannian distance. Let $R:G\to O(E)$ be a continuous orthogonal representation on a finite-dimensional real Euclidean space. Fix $z_0\in E$ of norm one, put
\[
 \mathcal O=R(G)z_0,\qquad Q=\operatorname{conv}(\mathcal O),
 \qquad q=\dim\mathcal O,
\]
and assume that $\mathcal O$ spans $E$ and that $E$ has no nonzero fixed vector. Thus $q\ge1$. A finite alphabet in $G$ contains the identity and carries a probability law $\mu$ for which
\begin{equation}\label{eq:groupgap57}
 \|P^nf\|_2\le\lambda^n\|f\|_2\quad\left(\int_Gf\,dm_G=0\right),
 \qquad Pf(g)=\E_{a\sim\mu}f(u_ag),\quad 0<\lambda<1.
\end{equation}
Letters outside the support of $\mu$ are allowed. All command products and stochastic rows retain the conventions of Definition~\ref{def:machine54}.

The uniform geometric hypothesis is that, for some $p>0$, $A<\infty$ and $h_0>0$,
\begin{equation}\label{eq:smallball57}
 \sup_{u,v\in S(E)}m_G\{g:\|R(g)u-v\|<h\}\le Ah^p
                      \qquad(0<h\le h_0).
\end{equation}
It is imposed on all unit vectors $u$, not only on $\mathcal O$. A conditional mean of points in $\mathcal O$, after normalization, need not belong to $\mathcal O$.

\begin{lemma}[From orbit small balls to executable contraction]\label{lem:orbitgap57}
Under \eqref{eq:groupgap57}--\eqref{eq:smallball57}, constants $c,C>0$ and integers $B_k\le C\log(k+2)$ can be chosen so that
\begin{equation}\label{eq:orbitgap57}
 \E_{w\sim\mu^{B_k}}\max_{1\le i\le k}
       \langle v_i,R(u_w)u\rangle\le1-c k^{-2/p}
\end{equation}
for every $u,v_1,\ldots,v_k\in S(E)$. The constant $c$ can be decreased so that the defect lies in $(0,1)$ for every $k\ge1$.
\end{lemma}
\begin{proof}
Choose $h=c_0k^{-1/p}\le\min(h_0,1)$ with $A c_0^p\le1/2$. Outside the union of the $k$ open balls in \eqref{eq:smallball57},
\[
 1-\max_i\langle v_i,R(g)u\rangle\ge h^2/2.
\]
The union has measure at most $1/2$. Consequently the Haar mean of
$f(g)=\max_i\langle v_i,R(g)u\rangle$ is at most $1-a_k$, where
$a_k=c_0^2k^{-2/p}/4$. Repeated centers are allowed. All such functions are bounded by one in absolute value and have a common Lipschitz constant $L$: differentiating the finite-dimensional representation bounds the derivative uniformly on unit vectors.

Write $D=\dim G$, distinct from the cap exponent $p$ and orbit dimension $q$. The ball calculation uses the fixed bi-invariant Riemannian metric on $G$. In exponential coordinates at the identity, Haar measure has a smooth positive density. Hence a metric ball of radius $\delta\le\delta_0$ has measure at least $c_G\delta^D$. Average $P^Bf$ over this ball. Left translations are isometries, so $P^Bf$ has Lipschitz constant at most $L$. Cauchy--Schwarz and \eqref{eq:groupgap57} give
\[
 |P^Bf(1)-\bar f|
 \le L\delta+2c_G^{-1/2}\delta^{-D/2}\lambda^B.
\]
Take $\delta$ to be a sufficiently small fixed multiple of $a_k$, also bounded by $\delta_0$. Then choose $B_k\ge1$ to make the second term at most $a_k/4$. The first term is at most $a_k/4$ as well. Since $\log(a_k^{-1})=O(\log(k+2))$, this requires $B_k\le C\log(k+2)$. Products under $\mu^{B_k}$ are genuine words of one common length. Their reverse chronological order has the same law because the letters are independent with the same law. Thus $P^{B_k}f(1)$ is the expectation in \eqref{eq:orbitgap57}. This proves the lemma with defect $a_k/2$.
\end{proof}

\begin{lemma}[Quadratic inner approximation of an orbit hull]\label{lem:orbithull57}
There are constants $C_1,C_2,\delta_0>0$ such that, for $0<\delta\le\delta_0$, points $z_1,\ldots,z_K\in\mathcal O$, including $z_0$, can be chosen with
\begin{equation}\label{eq:orbithull57}
 K\le C_1\delta^{-q},\qquad
 (1-C_2\delta^2)Q\subseteq\operatorname{conv}\{z_1,\ldots,z_K\}\subseteq Q,
 \qquad 0<1-C_2\delta^2<1.
\end{equation}
\end{lemma}
\begin{proof}
The orbit is a compact smooth embedded homogeneous manifold. Finitely many coordinate charts with bounded derivatives and inverse derivatives give an intrinsic $\delta$-net of size at most $C\delta^{-q}$; adjoining $z_0$ changes only $C$. Use the induced Riemannian metric on the orbit. If the functional $x\mapsto\langle b,x\rangle$ attains its maximum at $x_b\in\mathcal O$, its derivative along the orbit vanishes there. The second derivative along unit-speed orbit geodesics has absolute value at most $C\|b\|$, uniformly in their starting points and directions, by compactness of the unit tangent bundle and smoothness of the embedding. A net point within intrinsic distance $\delta$ of $x_b$ therefore satisfies
\[
 \langle b,z_i\rangle\ge h_Q(b)-C\|b\|\delta^2,
 \qquad h_Q(b)=\max_{x\in Q}\langle b,x\rangle.
\]
Haar averaging gives $\int R(g)z_0\,dm_G=0$. Moreover $0$ lies in the interior of $Q$ in $E$. Otherwise a nonzero supporting functional would be nonnegative on the orbit and have Haar mean zero; continuity and full Haar support would make it vanish on the orbit, contrary to its spanning $E$. Thus $r\mathbb B_E\subset Q$ for some $r>0$, and $h_Q(b)\ge r\|b\|$. The support inequality becomes
$\max_i\langle b,z_i\rangle\ge(1-C\delta^2/r)h_Q(b)$.
The support-function criterion for inclusion of compact convex sets proves \eqref{eq:orbithull57}. Decrease $\delta_0$ to obtain the last bounds.
\end{proof}

\begin{theorem}[Geometric width and occupation]\label{thm:uniformorbit57}
Assume \eqref{eq:groupgap57}--\eqref{eq:smallball57}. Take one seed and one query with target $F(g)=\rho R(g)z_0$, output set $Q$, and Euclidean error. For $0<\rho\le1$ and $0\le\varepsilon<\rho$, put $\zeta=\rho-\varepsilon\in(0,1]$. Every error-$\varepsilon$ realization satisfies
\begin{equation}\label{eq:orbitoccupation57}
 \#\{1\le t\le N:|S_t|\le k\}
 \le C\log(k+2)\left(1+k^{2/p}\log\frac1\zeta\right),
\end{equation}
where $C$ depends on the geometric and spectral-gap data but not on $\rho,\varepsilon,N,k$ or the realization. In particular, at fixed parameters,
\begin{equation}\label{eq:orbitrates57}
 W_{N,\varepsilon}\ge c\left(\frac{N}{\log(N+2)}\right)^{p/2}.
\end{equation}
For $0<\rho<1$ there is an exact common-row realization with
\begin{equation}\label{eq:orbitupper57}
 W_{N,0}\le C\left(1+\frac{N+1}{-\log\rho}\right)^{q/2}.
\end{equation}
For $\rho=1$ and fixed $0<\varepsilon<1$, the upper order is $O_\varepsilon((N+1)^{q/2})$. When $p=q$, the lower and upper exponents agree, with the displayed logarithmic loss in the lower bound. The constrained error threshold is exactly $\rho$, attained by the constant zero decoder.
\end{theorem}
\begin{proof}
Put $Z=R(u_w)z_0$ and let $D(S)\in Q\subseteq\mathbb B_E$ be the decoder. For each deterministic word, the error bound gives
$\langle\E[D(S)\mid w],Z\rangle\ge\rho-\varepsilon$.
Under any external word law, conditioning on the register yields
\[
 \zeta\le\E\langle D(S),Z\rangle
          \le\E\|\E[Z\mid S]\|.
\]
The conditional-centroid Lemma~\ref{lem:centroid54} applies to the law of Lemma~\ref{lem:orbitgap57}. The uniformity over every direction in \eqref{eq:smallball57} is precisely what justifies this application. It contracts the conditional norm by $1-a$ at a selected endpoint of width at most $k$, where $a=c k^{-2/p}\in(0,1)$. The initial norm is one. Fixed identity-letter fillers cannot increase it, regardless of their rows or intervening widths. After $J$ selected endpoints,
\[
 J[-\log(1-a)]\le\log(1/\zeta).
\]
Greedily select narrow cuts at separation at least $B_k$, omitting the first $B_k-1$ cuts. Each selected endpoint accounts for at most $B_k$ cuts. The count is at most $B_k(1+a^{-1}\log(1/\zeta))$, proving \eqref{eq:orbitoccupation57}, also when $\zeta=1$. In that case even one selected endpoint is impossible. More precisely one may take the integer bound
\[
 L_k=B_k\left(1+\left\lceil c^{-1}k^{2/p}\log(1/\zeta)\right\rceil\right).
\]
For a full width bound $k$, either $k\ge N^{p/2}$, or
$\log(k+2)\le C_p\log(N+2)$. In the latter case the occupation inequality
gives $N\le C_\zeta\log(N+2)k^{2/p}$. These alternatives prove
\eqref{eq:orbitrates57}.

For the upper bound take the net in Lemma~\ref{lem:orbithull57}, and write $r=1-C_2\delta^2$. The orbit hull is invariant under every command. Hence $rR(u_a)z_i$ is in the convex hull of the net points. For each of the finitely many command--vertex pairs, choose barycentric coordinates independently; no measurable selection is involved. Use them as row $i$ of a stochastic matrix $T_a$, the same at every epoch within this horizon. Initialize at $z_0$. The conditional row means then evolve as $r^tR(u_w)z_0$. Decode label $i$ by $(\rho/r^N)z_i$, which belongs to $Q$ if $r^N\ge\rho$, since $0\in Q$. Choosing
\[
 \delta^2=\min\left(\delta_0^2,\frac1{2C_2},
                          \frac{-\log\rho}{2C_2(N+1)}\right)
\]
gives this inequality by $\log(1-x)\ge-2x$ for $0\le x\le1/2$. The net size gives \eqref{eq:orbitupper57}. At unit amplitude and positive error use the exact construction for amplitude $1-\varepsilon/2$; its discrepancy from the desired target is $\varepsilon/2$.

Finally, Haar mean zero and $\|R(g)z_0\|=1$ imply
\[
 \int_G\|\rho R(g)z_0-y\|^2\,dm_G=\rho^2+\|y\|^2.
\]
Every constrained center therefore has worst error at least $\rho$, while $y=0$ attains it. Below it \eqref{eq:orbitrates57} rules out bounded width. None of the constructions requires transition matrices to satisfy the physical group relations, and extra finitely many commands do not change the upper construction.
\end{proof}
The geometric ingredients in Lemma~\ref{lem:orbithull57} are the usual local geometry of compact homogeneous spaces; related covering estimates are developed in \cite{Szarek57}. The additional realization statement is the common stochastic-row construction and its combination with a contraction uniform over conditional-centroid directions. A small-ball estimate restricted to the target orbit alone does not imply Theorem~\ref{thm:uniformorbit57}.

\begin{remark}[Dependence of constants]
The cap and executable-gap constants depend only on the fixed group with
its bi-invariant metric, the representation, the uniform cap data
$(A,p,h_0)$, and the declared $L^2$ gap $\lambda$. The hull constants also
depend on the target orbit, its second derivatives, and the inradius of
$Q$ in its span $E$. The occupation bound further depends on
$\zeta=\rho-\varepsilon$. None depends on the horizon, hidden transition
rows, or masses of individual labels. The orbit is the smooth embedded
quotient $G/G_{z_0}$ and $q$ is its real dimension. The logarithmic factor
is the loss of this mixing-block argument.
Theorem~\ref{thm:sharpoccupation61} removes it by a different
entropy--transport argument; all statements and proofs of the block
estimate remain valid.
\end{remark}
